\section{Identifiación del Problema}
Colombia históricamente ha tenido problemas con la ejecución de proyectos de 
infraestructura que están financiados por el Sistema General de Regalías (SGR) algunos 
de los problemas en los que se ha visto afectada han sido asociados con retrasos, 
sobrecostos, incumplimientos contractuales e incluso el abandono de obras, muchos de 
estos estan relacionados con la corrupción en la que está sumergido este pais, lo que ha 
generado los denominados “Elefantes Blancos” que son todos aquellos proyectos o 
infraestructura costosa de mantener que genera más gastos que beneficios, resultando 
que esta sea inútil y difícil de deshacerse de ella. Esta problemática no solo es generada 
por la ineficiencia en la gestión de los recursos. Sino que también genera impedimentos 
en el desarrollo social y económicos de las regiones beneficiaras. 

Pese a los esfuerzos institucionales para mejorar los mecanismos de control y vigilancia, 
La gran mayoría de las acciones de auditoria continúan siendo de carácter reactivo, es 
decir, toman acción una vez el daño se ha materializado o cuando se ve un contratiempo 
visible en la ejecución del proyecto. Lo anterior reduce la capacidad de las entidades de 
control para la prevención de irregularidades que puede haber en un proyecto público. 

En este entorno, la accesibilidad de grandes volúmenes de datos que provienen de 
plataformas como SECOP II, que es una plataforma transaccional en la cual las 
entidades estatales y los proveedores hacen toda la gestión en línea y en tiempo real  de 
todas las etapas de los procesos de contratación, en ellos encontramos procesos como: 
planeación, selección, ejecución y liquidación, Gracias a estas plataformas se abren 
nuevos caminos para el análisis predictivo a través de técnicas de Machine Learning. 
Pese a lo anterior en la actualidad no se cuenta con alguna herramienta sistemática que 
disponga de variables contractuales y financieras para hacer un predicción o 
anticipación desde la etapa de selección, por lo tanto, hay una alta probabilidad que los 
proyectos presenten contratiempos en su ejecución. 

Debido a esto, surge la necesidad de desarrollar una herramienta que implemente 
modelos predictivos que permita identificar de manera temprana contratos con alto 
riesgo de convertirse en “Elefantes Blancos” y proyectos que puedan estar bajo procesos 
irregulares, contribuyendo a que los entes de control cambien de un enfoque reactivo a 
uno preventivo. Este planteamiento no solo fortalecería la toma de decisiones de las 
entidades de control, sino que también promovería una gestión más eficiente, 
transparente y responsable en el manejo de los recursos públicos en proyectos estatales. 